\chapter{绪论}

\section{研究背景与意义}
流程工业广泛存在于化工、能源、冶金、制药、环保和材料制造等关键行业，其生产过程通常具有连续运行、变量耦合紧密、动态响应复杂、控制回路交织以及安全代价高昂等特点。一旦关键设备或工艺环节发生异常，不仅会导致产品质量波动、能耗升高和非计划停机，还可能进一步诱发连锁故障甚至安全事故。因此，如何及时、准确、稳定地识别工业过程中的故障，一直是过程监测、运行优化和安全保障中的核心问题~\cite{qin2012survey,ge2017review}。

随着分布式控制系统、工业互联网和边缘采集装置的普及，工业现场能够持续积累大规模、多变量、高频率的运行数据。与早期依赖机理模型、阈值规则和专家知识的诊断方法相比，数据驱动方法能够在较弱先验条件下从历史数据中自动提取异常模式，因此已经成为工业故障诊断的重要技术路线之一~\cite{taqvi2021review,yu2023challenges}。尤其是在深度学习方法进入故障诊断领域之后，一维卷积网络、循环网络和时序注意力模型显著增强了模型对复杂动态模式的表征能力，使得跨变量耦合和局部时序特征能够被更有效地学习~\cite{wu2018dcnn}。

然而，工业故障诊断面临的一个长期难题在于训练分布与部署分布往往并不一致。设备老化、原料组成变化、环境扰动、控制参数重整定、负载切换以及生产阶段变更等因素，都会造成不同时间段、不同工况乃至不同装置上的过程数据分布发生偏移。设训练阶段的数据分布为 $P_s(x,y)$，部署阶段的数据分布为 $P_t(x,y)$，在很多真实工业场景下有
\begin{equation}
P_s(x,y)\neq P_t(x,y).
\end{equation}
当模型仅依赖源域标注样本训练时，即使其在源域上取得很高准确率，也可能在目标域工况中出现显著性能下降。这一问题在流程工业中尤为突出，因为目标域故障样本通常稀缺、获取成本高且伴随安全风险，重新采集并人工标注大量目标域故障数据通常并不可行~\cite{pan2009survey,wu2020transfer}。

领域自适应作为迁移学习的重要分支，正是为缓解上述分布偏移问题而提出。其核心思想是在源域有标签、目标域少标签甚至无标签的条件下，通过统计分布对齐、对抗学习、类别结构约束或样本匹配等方式提升模型在目标域上的泛化能力。对于流程工业而言，如果能够在不依赖目标域大量标注样本的前提下稳定提升跨工况故障识别性能，将对缩短模型部署周期、降低迁移成本和增强工程可用性具有直接意义~\cite{wang2020lda,qin2023dynamic}。

Tennessee Eastman Process（TEP）是过程系统工程领域广泛使用的经典工业过程仿真平台，由 Downs 和 Vogel 提出~\cite{downs1993plant}。围绕 TEP 构造的领域自适应数据集将不同运行模式组织为不同领域，使得“跨工况故障诊断”能够被清晰表述为无监督领域自适应问题。该数据集既保留了工业多变量时序数据的动态特征，又提供了相对规范的域划分和故障标签，因此非常适合用于构建统一实验基准、比较典型方法并验证改进方法的有效性~\cite{montesuma2023benchmark}。

\section{国内外研究现状}
\subsection{工业过程故障诊断研究进展}
工业过程故障诊断的早期研究较多依赖过程机理、专家经验和手工规则。典型思路包括基于观测器或残差生成的模型方法、故障树分析、专家系统以及软测量建模等。这类方法在对象机理明确、系统结构相对稳定的场景下具有较好的解释性，但对大型流程工业而言，完整建立高精度机理模型并不容易；同时，实际运行中设备磨损、原料波动和控制策略调整会不断改变系统行为，使得纯机理方法在长期应用中面临维护成本高和泛化性不足的问题。

随着工业信息化水平提高，数据驱动方法逐渐成为流程工业故障诊断的主流研究方向。主成分分析、偏最小二乘、动态 PCA 和动态 PLS 等统计过程监测方法长期是流程工业故障检测的重要工具；支持向量机、近邻分类器、决策树和集成学习等浅层机器学习方法也被广泛用于故障分类任务~\cite{qin2012survey,ge2017review}。近年来，深度学习方法通过多层非线性映射自动学习故障相关表征，一维卷积网络能够从多变量时序窗口中提取局部动态模式，循环网络与注意力结构则在捕捉长时依赖关系方面具有优势~\cite{wu2018dcnn,taqvi2021review}。然而，当工况切换导致训练集与测试集分布不一致时，仅依赖源域监督学习得到的深度模型仍会暴露出显著的跨域泛化不足。

\subsection{迁移学习与领域自适应研究进展}
迁移学习旨在利用已有任务或已有域中的知识提升目标任务的学习效果。Pan 和 Yang 将迁移学习系统地定义为：在源域与目标域分布或任务不同的情况下，通过知识迁移改善目标任务学习性能的过程~\cite{pan2009survey}。当源域与目标域任务相同但数据分布不同，即
\begin{equation}
\mathcal{Y}_s=\mathcal{Y}_t,\qquad P_s(x)\neq P_t(x),
\end{equation}
便得到领域自适应问题。按照目标域是否具有标签，领域自适应可分为监督式、半监督式和无监督式等多种设置，本文重点关注目标域无标签的无监督领域自适应。

现有无监督领域自适应方法大体可以分为四类。第一类是统计分布对齐方法，例如 CORAL 通过协方差匹配减小二阶统计差异~\cite{sun2016deep}，DAN 通过最大均值差异在核空间中逼近边缘分布对齐~\cite{long2015dan}。第二类是对抗式方法，例如 DANN 通过域判别器和梯度反转层学习域不变特征~\cite{ganin2016dann}，CDAN 进一步将类别预测引入条件对齐过程，从而增强对类条件信息的感知~\cite{long2018conditional}。第三类是最优传输方法，例如 DeepJDOT 联合考虑特征距离与类别预测代价，在样本级别建立跨域匹配关系~\cite{damodaran2018deepjdot}。第四类则更强调目标域伪监督与类别结构，例如基于伪标签、一致性训练、类原型和聚类约束的自训练方法~\cite{lee2013pseudolabel,tarvainen2017meanteacher,tanwisuth2021prototype}。

总体来看，上述方法已经为跨工况诊断提供了坚实基础，但针对流程工业多类别时序任务仍存在三个明显缺口。其一，很多方法更偏向削弱整体域差异，而对目标域无标签样本如何稳定进入训练闭环讨论不足；其二，伪标签方法往往依赖单一置信度阈值，容易受到高置信错误和类别不均衡影响；其三，全局对齐虽然能够缩小源域和目标域的整体距离，却不能自动保证跨域同类样本聚合、异类样本分离，因而在多故障类别任务中容易出现类别混叠。

\subsection{TEP 跨工况迁移诊断研究现状}
Tennessee Eastman Process 由于过程结构完整、变量丰富、耦合复杂且具有较强代表性，长期被用于过程监测、故障检测和故障诊断算法评估~\cite{downs1993plant}。近年的研究开始把 TEP 扩展为跨域场景，将不同操作模式组织为不同领域，从而自然构造“跨工况迁移诊断”任务~\cite{montesuma2023benchmark,wu2020transfer,qin2023dynamic}。相关工作已经表明，不同模式之间的分布偏移会显著影响模型迁移效果，而统一任务定义、数据划分、归一化方式、骨干网络结构和模型选择策略，对于公平比较不同领域自适应方法至关重要。

已有研究在 TEP 跨工况迁移诊断中取得了不少进展，但仍常见以下问题：不同工作采用的数据划分和训练协议不完全一致，导致横向比较困难；部分方法只关注少量场景，难以判断其是否具有跨场景稳定性；一些改进方法虽然能在个别任务上取得收益，但对“为什么有效、在哪些场景有效、哪些模块真正发挥作用”的分析仍然不足。基于这些现实缺口，本文将研究重点放在两条主线上：一是构建面向当前工作区实现的统一 benchmark 组织方式，二是围绕多类别工业时序迁移中的稳定伪监督与类条件结构建模提出任务导向的改进方法。

\section{研究内容与技术路线}
围绕跨工况无监督领域自适应故障诊断问题，本文的研究内容主要包括以下五个方面。

\begin{enumerate}
\item 面向 TEP 多工况迁移诊断任务，建立统一的数据读取、预处理、训练、评估、结果导出与图表汇总流程，在共享协议下组织领域自适应实验；
\item 基于固定折划分构建包含 11 个迁移场景、8 种方法的 benchmark，其中既包含 8 个单源到单目标场景，也包含 3 个多源到单目标场景，从而兼顾横向比较广度与训练成本可控性；
\item 在统一 FCN 骨干网络和统一结果产物结构下，比较 Source Only、Target Only、CORAL、DAN、DANN、CDAN、DeepJDOT 与 RCTA 的表现，形成可复现的实验基线；
\item 针对基线比较中暴露出的目标域伪标签噪声、类别混叠以及训练阶段不稳定问题，提出可靠性门控的类条件时序自适应方法 RCTA，并将其拆解为教师--学生时序一致性、可靠性门控、双原型记忆和课程化优化四个递进章节；
\item 在实验章节中，将“当前已需完成的 88-run benchmark”和“后续模块级消融计划”加以区分：前者用于形成主结果，后者用于在有额外算力预算时进一步验证各模块的独立贡献。
\end{enumerate}

本文的技术路线可以概括为如下四步。首先，结合 TEP 领域自适应数据集明确域定义、输入表示、归一化方式、模型选择准则和固定折策略，为后续比较设定统一边界。其次，在共享骨干和统一训练脚手架下复现代表性领域自适应方法，观察不同对齐机制在 11 个 benchmark 场景中的表现。再次，根据基线方法的共性问题，梳理“目标域样本如何稳定进入训练”“哪些目标样本可信”“如何保持类条件结构”“多项约束如何按节奏进入优化”四个关键问题。最后，围绕上述问题设计 RCTA 的四个模块，并在统一 benchmark 上验证完整方法的收益，同时为后续严格消融预留结构化接口。

\section{研究问题、核心挑战与创新点}
基于前述研究现状，本文关注的核心问题是：在目标域缺少故障标签的条件下，如何利用源域标注样本与目标域无标注样本，训练一个能够跨工况迁移的故障诊断模型。本文的方法推导以单源到单目标无监督领域自适应为主线，同时在实验上引入少量多源到单目标场景，以观察方法在更复杂源域组合下的迁移稳定性。

为了使创新点与问题痛点形成一一对应，本文将核心挑战与后文创新章节组织如下。

\begin{enumerate}
\item \textbf{挑战一：跨工况全局分布偏移显著，但单纯对齐不足以稳定利用目标域时序样本。} 不同运行模式下的负载水平、控制策略和动态响应不同，目标域样本虽然可见但无标签，模型难以直接获得目标域分类监督。针对这一痛点，第三章在强对齐底座上引入 EMA 教师网络和弱强双视图增强，使目标域样本以更平滑的方式进入训练闭环。
\item \textbf{挑战二：目标域伪标签容易受噪声和类别不均衡影响。} 在 29 类故障诊断任务中，若直接使用高置信伪标签，容易放大早期错误并让少数容易类别占据伪监督预算。针对这一痛点，第四章提出多因素可靠性评分、按类门控与样本分层利用机制，尽可能把目标域监督集中到更可信的样本上。
\item \textbf{挑战三：全局域对齐容易带来类别混叠，难以保证跨域同类样本聚合、异类样本分离。} 对于类间边界复杂的工业时序任务，简单缩小整体域差异可能引发负迁移。针对这一痛点，第五章提出双原型记忆驱动的类条件结构约束，通过源域原型、目标域原型和参考原型共同塑造几何结构。
\item \textbf{挑战四：多个辅助目标的作用时机不同，若同时强行注入可能破坏训练稳定性。} 伪标签、一致性和原型约束都依赖模型阶段性质量，过早或过强施加会干扰基础分类边界。针对这一痛点，第六章提出课程化优化与置信保护型模型选择策略，使各模块按“先搭底座、再引入目标域监督、最后强化类条件结构”的顺序逐步发挥作用。
\end{enumerate}

\subsection{创新定位说明}
需要强调的是，本文的创新并非宣称在领域自适应理论上提出了完全全新的基础范式，而是属于\textbf{面向工业多类别时序迁移任务的机制性集成创新}。其价值主要体现在以下三个层面。

\begin{enumerate}
\item \textbf{问题驱动的模块组合。} 本文不是简单堆叠已有技巧，而是围绕目标域样本引入、伪标签筛选、类条件结构约束和训练节奏控制四个具体失效点，构造出有前后依赖关系的完整链条。
\item \textbf{实现对齐与可复现比较。} 所有方法均在同一套骨干网络、训练脚手架、固定折策略和结果导出结构下运行，使得方法比较尽量聚焦于自适应机制本身，而非实现细节差异。
\item \textbf{面向工程任务的边界控制。} 本文在写作与设计上刻意避免夸大“单一模块的理论创新”，而更强调“任务痛点识别准确、方法设计闭环清晰、实验协议可复现、改进收益可验证”这一类本科论文更稳健的贡献形式。
\end{enumerate}

据此，本文的主要创新点可概括为：
\begin{enumerate}
\item 构建面向当前工作实现的 TEP 固定折 benchmark，将 11 个跨工况场景和 8 种方法纳入同一实验协议，提供可复现的横向比较基础；
\item 提出教师--学生时序一致性机制，通过 EMA 教师与弱强双视图增强，为目标域样本提供更平滑的伪监督入口；
\item 提出多因素可靠性门控和分层利用策略，从置信、熵、一致性和原型相似度四个维度联合评估伪标签质量；
\item 提出双原型记忆驱动的类条件结构约束，用参考原型、吸引损失和分离损失缓解类别混叠；
\item 提出课程化优化与置信保护型选择策略，将门控课程、损失预热和 checkpoint 选择统一纳入完整训练流程。
\end{enumerate}

\section{论文结构安排}
全文共分为八章，各章内容安排如下。

第一章为绪论，介绍研究背景、国内外研究现状、研究内容、核心挑战以及本文创新点，并明确本文的创新定位。第二章为跨工况迁移诊断的基准脉络，说明任务定义、统一协议、共享骨干网络、典型基线方法和模型选择原则。第三章为教师--学生时序一致性自适应，阐述 RCTA 的总体框架、EMA 教师网络以及弱强双视图一致性机制。第四章为可靠性门控的伪标签学习，说明多因素可靠性评分、按类门控和样本分层利用方式。第五章为双原型记忆的类条件结构约束，介绍源域/目标域原型记忆、参考原型构造以及吸引与分离机制。第六章为课程化优化与完整训练流程，给出门控课程、损失预热、模型选择策略和算法流程。第七章为实验设计与结果组织，围绕当前需要执行的 88-run fixed-fold benchmark 展开，并预留后续模块消融位置。第八章为总结与展望，对全文工作进行归纳，并讨论后续更强创新方向的升级路径。
